\section{Threats to validity, alternative explanations and decisive falsifiers}
\label{sec:threats}

A paper about self-improving research systems is unusually vulnerable to circular evidence. We therefore state the strongest alternative explanations and the outcomes that would force claim narrowing.

\subsection{The live case study is observational and selected}

The six Millennium programmes were intentionally chosen because they are difficult, open and structurally diverse. They are not a random sample of scientific research. Their value is as a stress-test observatory, not as a population from which we can directly estimate the average behavior of all research agents. Furthermore, individual atomic tasks are selected adaptively by the agents and framework. Frequencies observed in the live archive therefore mix properties of the problem distribution with properties of the selection policy.

Paper 5 accordingly separates retrospective pathology counts from prospective benchmark results. Statements such as ``surrogate/root-coordinate confusion recurs across several domains'' are case-study observations. Population-level rates require a frozen sampling rule or benchmark set.

\subsection{Open problems do not provide frequent final success labels}

For an unsolved root problem there may be no positive terminal outcome during the study. This creates a risk that the framework rewards internally generated notions of progress. Residual contraction, lesson authority and local theorem validation are therefore reported as intermediate coordinates, never as substitutes for root success. A local lemma can be mathematically correct while contributing nothing to the root. The gluing/interface metrics explicitly measure this failure mode.

The strongest capability claims will therefore rely on auxiliary task distributions with externally checkable or held-out outcomes in addition to the live open-problem observatory.

\subsection{The underlying model can be the real cause}

An agent may solve a task because the language model already knew how. This is the motivation for the four-arm design in Section~\ref{sec:metrology}. If \texttt{MODEL\_ONLY} performs as well as \texttt{RAKL\_LEARNING} at matched validity and resources, the data do not support a Orion capability benefit. If learning and sham-memory arms are similar, the content-specific memory hypothesis is not supported. If \texttt{RAKL\_RESET} accounts for the gain, the effect belongs to static workflow scaffolding rather than continual experience.

\subsection{Extra compute and context can masquerade as intelligence}

Orion performs retrieval, preprocessing, validation and memory compilation. An unbudgeted comparison against a shorter direct prompt would be unfair. All arms therefore share a preregistered resource ceiling, and actual usage is reported. The sham-memory arm additionally controls, imperfectly, for the possibility that merely adding more tokens or structured context improves performance.

No sham is perfect. Irrelevant memory may distract more than an empty context, while carefully chosen near-miss memory may create a harder condition. We will therefore freeze sham construction rules before task outcomes and report more than one sham stratum when sufficient budget is available.

\subsection{Prompt and workflow differences are part of the intervention}

A full Orion-versus-model-only comparison necessarily changes instructions. It estimates the effect of the complete architecture, not a single isolated variable. The reset-versus-learning comparison is cleaner for experience because both arms use the same Orion workflow. Component ablations then localize specific mechanisms. We will not interpret the total Orion lift as a pure memory effect.

\subsection{Model-service drift and non-determinism}

Hosted model names can conceal backend changes, and repeated samples can differ even under nominally fixed settings. Exact model revision is recorded when available. Strong effects should be replicated across seeds/sampling configurations and, where the claim is intended to generalize, across more than one model family or revision. Failure to reproduce under a changed hosted model does not by itself distinguish framework failure from model drift, but it weakens claims of model-independent benefit.

\subsection{Temporal dependence and pseudo-replication}

Development episodes are not independent because later states contain earlier experience. Treating every cycle as an IID sample would underestimate uncertainty. The task is the unit for paired benchmark comparisons; fresh-transfer tasks all start from the same frozen learned state. Longitudinal curves remain descriptive unless the dependence structure is explicitly modelled. Multiple episodes generated from the same source or failure lineage are not counted as independent transfer validations.

\subsection{Concurrent agents can contaminate lanes}

Several research agents write to a shared application repository. One lane can therefore encounter artifacts produced by another lane between its start and completion. Exact base/head identities and fibre snapshots are required to establish what was available at decision time. A fresh-transfer benchmark uses frozen immutable states rather than a moving shared \texttt{main}. Uncontrolled cross-lane writes invalidate a benchmark packet but remain legitimate observational case evidence.

\subsection{Telemetry changes the behavior being observed}

Requiring agents to record failures, retrieved memory and saturation may itself improve or constrain behavior. This is not a defect if the telemetry is part of Orion, but it matters for interpretation. The pre-telemetry archive and reset/model-only arms provide partial controls. If Paper 5 claims that observation alone reveals natural model behavior, that claim would be too strong. We study the behavior of agents operating inside a specified research architecture.

\subsection{Semantic novelty labels can be wrong}

The seven-axis growth vector requires identity and semantic-retention judgments. Raw counts are objective relative to the stored data; retained novelty can depend on imperfect canonicalization. We therefore preserve the raw object stream, version novelty criteria, permit later supersession and avoid using retained novelty as a sole promotion metric. Periodic external/manual audits of sampled novelty assignments are desirable.

\subsection{Residual contraction can be gamed}

A framework can make the blocker set smaller by hiding or coarsening obligations. This would create artificial progress. Residual transformations therefore preserve evidence and are cross-checked against root contracts and gluing coverage. A step that exposes new independent blockers is not penalized merely because the count increases. Hostile tests should include obstruction renaming and deliberate omission of an interface.

\subsection{Retrieval recall requires a bound universe}

A claim that no relevant memory was found is not measurable against an unbounded repository narrative. The XM003 case exposed exactly this problem. Retrieval recall is therefore reported only relative to a content-identified index/universe and later relevance adjudication. A future framework-level coverage receipt is an explicit engineering requirement rather than an assumed property.

\subsection{Same-session expert cells are not independent reviewers}

The live research lanes use role-separated expert passes for adversarial coverage. These roles share model context and cannot be counted as independent peer review. The Paper 5 drafting/review loop in the present session has the same limitation. Nature-style concern taxonomies can improve coverage, but same-session internal review is labelled as such. Strong external-review claims require genuinely separate people/processes and evidence lineage where relevant.

\subsection{Researcher and author intervention}

The human operator chooses high-level goals, can change scheduled-agent prompts, may override merges and is co-designing the framework while observing it. The system is therefore not an experiment in fully autonomous AI development. Operator interventions are timestamped as part of the method history. Paper 5's claim concerns evidence-governed human/AI recursive research infrastructure, not human-free self-evolution.

\subsection{The paper is being written while the system evolves}

A live repository creates a moving-target risk. Every reported result therefore needs an explicit evidence cutoff and exact subject. Later findings can be added as new versions of the manuscript, but they cannot be silently backdated into the original evaluation. The present draft intentionally contains placeholders for prospective metrics rather than estimating them from incomplete future data.

\subsection{Publication and confirmation bias}

A self-improvement project naturally prefers stories of successful evolution. We counter this by preserving rejected variants, baseline-only wins, invalid transfers, failed consolidations, blocked promotions and operational overrides. Paper 5 should publish a complete registered evaluation packet, not only the tasks on which Orion helps. If the result is null or negative, the framework and paper claims must narrow accordingly.

\subsection{Security and evaluator independence are assumptions, not proofs}

Content addressing and protected paths do not protect against a compromised repository administrator, CI service, dependency or assurance operator. The no-self-promotion proposition is conditional on these boundaries. Threats outside the trusted computing base are reported rather than treated as solved by hashes.

\subsection{Decisive result branches}

The prospective programme is designed so that important hypotheses can fail cleanly.

\paragraph{H1: accumulated Orion experience improves fresh research.}
Supported only if \texttt{RAKL\_LEARNING} shows a material preregistered gain over \texttt{RAKL\_RESET} on fresh transfer without a blocking validity/resource regression. Refuted for the tested distribution if the effect is materially negative. Inconclusive if confidence/coverage is insufficient.

\paragraph{H2: learned memory content, not extra context, causes part of the gain.}
Supported only if learning materially outperforms preregistered sham-memory controls. If sham and learned memory are equivalent, the content-specific claim is unsupported even if both outperform reset.

\paragraph{H3: Orion helps beyond the base LLM.}
Supported in scope only if total matched lift over \texttt{MODEL\_ONLY} is positive on the registered primary QoIs or if a separately preregistered epistemic-safety benefit is material. A substantial excess of \texttt{BASELINE\_ONLY\_SUCCESS} tasks is evidence against the claim.

\paragraph{H4: experience reduces repeated structural failure.}
Supported if fresh-task repeated-failure rate falls without increasing false-transfer or invalid-authority rates. If the framework simply avoids old failures by making different but equally serious failures, the hypothesis is not supported.

\paragraph{H5: framework evolution is more than benchmark overfit.}
Supported for a particular parent--child edge only after a preregistered development gain transfers to a fresh, protected, lineage-independent assurance packet under matched resources. Development-only improvement is explicitly classified as local improvement rather than evolution evidence.

\paragraph{H6: Orion's operator/representation atlas approaches a high zero-invention regime.}
This Orion-triviality hypothesis remains open. A growing share of verified tasks solved with only preexisting resources would support it; a persistent need for new representations/operators/ontology across fresh tasks would refute or bound it.

These result branches are part of the method. They prevent future manuscript revisions from defining ``success'' after observing whichever metrics happen to improve.